%% kapa: exam=1 round=35 year=2024 period=1 form=A subject=economics subject_ko=경제학원론 seq=32 number=72
%% source: Q-Net 감정평가사 기출문제 (cst00309) artlSeq=5213229 file=raw/1008/5213229/1교시 문제지 원본.pdf
%% conversion: kordoc 4.24.51 -> tools/gen1.py (auto); review=proofread-sonnet; flags=none
\begin{question}{72}
한국과 미국의 인플레이션율이 각각 3 \%와 5 \%이다. 구매력평가설과 이자율평가설(interest parity theory)이 성립할 때, 미국의 명목이자율이 5 \%라면, 한국의 명목이자율은? (단, 기대인플레이션율은 인플레이션율과 동일하다.)
\begin{choices}
  \item 1 \%
  \item 2 \%
  \item 3 \%
  \item 4 \%
  \item 5 \%
\end{choices}
\end{question}
\answer{3}
